\documentclass[10pt,a4paper]{amsart}
\usepackage{amsmath,amssymb,amsthm,graphicx}
\usepackage{verbatim}
\newenvironment{reference}{\comment}{\endcomment}
%\newenvironment{reference}{}{}
\newenvironment{slogan}{\comment}{\endcomment}
\newenvironment{history}{\comment}{\endcomment}

% For commutative diagrams we use Xy-pic
\usepackage[all]{xy}

% We use 2cell for 2-commutative diagrams.
\xyoption{2cell}
% Theorem environments.
%
\theoremstyle{plain}
\newtheorem{theorem}[subsection]{Theorem}
\newtheorem{proposition}[subsection]{Proposition}
\newtheorem{lemma}[subsection]{Lemma}

\theoremstyle{definition}
\newtheorem{definition}[subsection]{Definition}
\newtheorem{example}[subsection]{Example}
\newtheorem{exercise}[subsection]{Exercise}
\newtheorem{situation}[subsection]{Situation}

\theoremstyle{remark}
\newtheorem{remark}[subsection]{Remark}
\newtheorem{remarks}[subsection]{Remarks}

\numberwithin{equation}{subsection}

% Macros
%
\def\lim{\mathop{\mathrm{lim}}\nolimits}
\def\colim{\mathop{\mathrm{colim}}\nolimits}
\def\Spec{\mathop{\mathrm{Spec}}}
\def\Hom{\mathop{\mathrm{Hom}}\nolimits}
\def\Ext{\mathop{\mathrm{Ext}}\nolimits}
\def\SheafHom{\mathop{\mathcal{H}\!\mathit{om}}\nolimits}
\def\SheafExt{\mathop{\mathcal{E}\!\mathit{xt}}\nolimits}
\def\Sch{\mathit{Sch}}
\def\Mor{\mathop{\mathrm{Mor}}\nolimits}
\def\Ob{\mathop{\mathrm{Ob}}\nolimits}
\def\Sh{\mathop{\mathit{Sh}}\nolimits}
\def\NL{\mathop{N\!L}\nolimits}
\def\CH{\mathop{\mathrm{CH}}\nolimits}
\def\proetale{{pro\text{-}\acute{e}tale}}
\def\etale{{\acute{e}tale}}
\def\QCoh{\mathit{QCoh}}
\def\Ker{\mathop{\mathrm{Ker}}}
\def\Im{\mathop{\mathrm{Im}}}
\def\Coker{\mathop{\mathrm{Coker}}}
\def\Coim{\mathop{\mathrm{Coim}}}

% Boxtimes
%
\DeclareMathSymbol{\boxtimes}{\mathbin}{AMSa}{"02}

%
% Macros for moduli stacks/spaces
%
\def\QCohstack{\mathcal{QC}\!\mathit{oh}}
\def\Cohstack{\mathcal{C}\!\mathit{oh}}
\def\Spacesstack{\mathcal{S}\!\mathit{paces}}
\def\Quotfunctor{\mathrm{Quot}}
\def\Hilbfunctor{\mathrm{Hilb}}
\def\Curvesstack{\mathcal{C}\!\mathit{urves}}
\def\Polarizedstack{\mathcal{P}\!\mathit{olarized}}
\def\Complexesstack{\mathcal{C}\!\mathit{omplexes}}
% \Pic is the operator that assigns to X its picard group, usage \Pic(X)
% \Picardstack_{X/B} denotes the Picard stack of X over B
% \Picardfunctor_{X/B} denotes the Picard functor of X over B
\def\Pic{\mathop{\mathrm{Pic}}\nolimits}
\def\Picardstack{\mathcal{P}\!\mathit{ic}}
\def\Picardfunctor{\mathrm{Pic}}
\def\Deformationcategory{\mathcal{D}\!\mathit{ef}}
\title{Extension of algebraically closed fields does not change the kernel of the Frobenius linear map minus one}
\author{The Stacks Project Authors}
\date{stacks-project:etale-cohomology.tex:local-snapshot-20261011; GNU FDL 1.2 or later}
\begin{document}
\maketitle
\noindent\textit{Editorial note.} All delivered mathematical text is the author's own, copied from the pinned Stacks Project chapter file etale-cohomology.tex; no statement or proof step was rewritten, reordered or removed.\par\medskip
\noindent\textit{Editorial note.} The proof is the author's complete argument: the coordinate description of the kernel is written out, the finiteness input is attributed to the companion lemma, and the final base change step is stated.\par\medskip
\noindent\textit{Editorial note.} Every labelled reference occurring inside the counted spans is carried in assembly.external\_references as a non-load-bearing navigation pointer (same-chapter ones at their definition site, cross-chapter ones at the invocation inside the counted span), so no load-bearing mathematics is replaced by a stub. No bibliography citation occurs inside the counted spans.\par\medskip
\noindent\textit{Editorial note.} The checked chapter list for same-name collisions is the set of stacks chapters staged in this repository: chow, cohomology, divisors, etale-cohomology, flat, more-algebra, more-morphisms, morphisms, spaces-more-morphisms, varieties together with this card's chapter etale-cohomology.\par\medskip
\noindent\textit{Editorial note.} The preamble is the chapter's own native preamble excerpt already used for this source: the Xy-pic setup (native preamble lines 14-18) and the native theorem declarations with the mathematical macros (native preamble lines 154-217).\par\medskip
\noindent\textit{Allowed standard prerequisites (editorial).}\begin{itemize}
\item The situation and conclusion of Lemma F-1. Same chapter, Lemma lemma-F-1: in the standing situation the equations cutting out the kernel define a closed subscheme of affine space finite over Spec(k). Supplies both the situation the statement refers to and the finiteness used in the proof.
\item Points of a finite scheme over an algebraically closed field after an extension. Standard base change behaviour of the k-valued points of a scheme finite over Spec(k) when k is algebraically closed, as used by the author in the closing sentence. Supplies the equality Z(k) = Z(k') that finishes the proof.
\end{itemize}\medskip
\begin{lemma}
\label{lemma-F-1}
Let $k$ be an algebraically closed field of characteristic $p > 0$.
Let $V$ be a finite dimensional $k$-vector space. Let $F : V \to V$
be a frobenius linear map, i.e., an additive map such that
$F(\lambda v) = \lambda^p F(v)$ for all $\lambda \in k$ and $v \in V$.
Then $F - 1 : V \to V$ is surjective with kernel a finite dimensional
$\mathbf{F}_p$-vector space of dimension $\leq \dim_k(V)$.
\end{lemma}

\begin{lemma}
\label{lemma-F-2}
In the situation of Lemma \ref{lemma-F-1} let $k'/k$ be an extension
of algebraically closed fields. Set $V' = V \otimes_k k'$ and
$F' : V' \to V'$ be the Frobenius linear map induced by $F$.
Then $\Ker(F - 1)$ maps isomorphically onto $\Ker(F' - 1)$.
\end{lemma}

\begin{proof}
Let $v_1, \ldots, v_n$ be a basis of $V$ and say $F(v_i) = \sum a_{ij}v_j$.
Then $v = \sum \lambda_i v_i$ is in $\Ker(F - 1)$ if and only if
$\sum \lambda_i^pa_{ij} = \lambda_j$ for $j = 1, \ldots, n$.
By Lemma \ref{lemma-F-1} these equations
cut out a closed subscheme $Z$ of
$\mathbf{A}^n_k = \Spec(k[\lambda_1, \ldots, \lambda_n])$
finite over $\Spec(k)$. Then $Z(k) = Z(k')$ and we conclude.
\end{proof}

\end{document}
